\documentclass[a4paper, serif, 10pt]{moderncv}

%% moderncv
% color and style
\moderncvcolor{black}
\moderncvstyle{banking}

% renew moderncv command to adapt for CJK colon
\renewcommand*{\cvitem}[3][.25em]{%
  \ifstrempty{#2}{}{\hintstyle{#2}：}{#3}%
  \par\addvspace{#1}}

\renewcommand*{\cvitemwithcomment}[4][.25em]{%
  \savebox{\cvitemwithcommentmainbox}{\ifstrempty{#2}{}{\hintstyle{#2}：}#3}%
  \setlength{\cvitemwithcommentmainlength}{\widthof{\usebox{\cvitemwithcommentmainbox}}}%
  \setlength{\cvitemwithcommentcommentlength}{\maincolumnwidth-\separatorcolumnwidth-\cvitemwithcommentmainlength}%
  \begin{minipage}[t]{\cvitemwithcommentmainlength}\usebox{\cvitemwithcommentmainbox}\end{minipage}%
  \hfill% fill of \separatorcolumnwidth
  \begin{minipage}[t]{\cvitemwithcommentcommentlength}\raggedleft\small\itshape#4\end{minipage}%
  \par\addvspace{#1}}

%% page layout/margins
\usepackage[top=2.5cm, bottom=2.5cm, left=1.5cm, right=1.5cm]{geometry}

\nopagenumbers{}

%% Babel config for English language
\usepackage[english]{babel}

%% fontspec
\usepackage{fontspec}

\IfFontExistsTF{Linux Libertine}{
  \setmainfont[Ligatures={TeX, Common}, Numbers=Lining, ItalicFont=Linux Libertine]{Linux Libertine}
}{}
\IfFontExistsTF{Linux Libertine O}{
  \setmainfont[Ligatures={TeX, Common}, Numbers=Lining, ItalicFont=Linux Libertine O]{Linux Libertine O}
}{}

%% CTeX
% CJK support, used to show CJK characters in the resume
%
% - heading=false: disable ctex heading
% - punct=kaiming: use kaiming punctuations styles for CJK
% - scheme=plain: use plain scheme, do not override \normalsize font size
% - space=auto: space settings for CJK characters
%
% CJK fonts are resolved per locale language, e.g. Noto Serif/Sans CJK JP
% for Japanese resumes, with fallback to the next region if not installed
%
% ref:
% - http://ctan.mirrorcatalogs.com/language/chinese/ctex/ctex.pdf
\usepackage[UTF8, heading=false, punct=kaiming, scheme=plain, space=auto]{ctex}

\IfFontExistsTF{Noto Serif CJK HK}{
  \setCJKmainfont{Noto Serif CJK HK}
}{
\IfFontExistsTF{Noto Serif CJK TC}{
  \setCJKmainfont{Noto Serif CJK TC}
}{}}
\IfFontExistsTF{Noto Sans CJK HK}{
  \setCJKsansfont{Noto Sans CJK HK}
}{
\IfFontExistsTF{Noto Sans CJK TC}{
  \setCJKsansfont{Noto Sans CJK TC}
}{}}

%% line spacing
\usepackage{setspace}
\setstretch{1.125}

%% URL styling - use normal text instead of monospace
\urlstyle{same}

%% Patch \httplink and \httpslink to handle full URLs
% The original moderncv macros always prepend http:// or https:// to URLs.
% This causes issues when the URL already has a protocol (e.g., https://example.com)
% resulting in malformed URLs like https://https://example.com.
% This patch checks if the URL already contains "://" and uses it directly if so.
% Note: We use \str_set:Nx (x-expansion) to expand macros like \@homepage first.
\ExplSyntaxOn
\str_new:N \g_yamlresume_url_str

\RenewDocumentCommand{\httpslink}{O{}m}{
  \str_set:Nx \g_yamlresume_url_str {#2}
  \str_if_in:NnTF \g_yamlresume_url_str {://}
    { \url{#2} }
    { \url{https://#2} }
}

\RenewDocumentCommand{\httplink}{O{}m}{
  \str_set:Nx \g_yamlresume_url_str {#2}
  \str_if_in:NnTF \g_yamlresume_url_str {://}
    { \url{#2} }
    { \url{http://#2} }
}
\ExplSyntaxOff

\name{陳家豪}{}
\title{資深網絡安全工程師 | 雲端安全與滲透測試}
\phone[mobile]{+852 9123 4567}
\email{ka.ho.chan@example.com}
\homepage{https://www.chanhk-security.com}

\address{中國香港，九龍，香港，香港九龍旺角彌敦道 700 號}{}{}

\extrainfo{{\small \faLinkedin}\ \href{https://www.linkedin.com/in/chankaho-cyber}{@chankaho-cyber} {} {} {} • {} {} {} 
{\small \faGithub}\ \href{https://github.com/chankaho-sec}{@chankaho-sec}}

\begin{document}

\maketitle

\section{簡介}

\cvline{}{具 8 年網絡安全工程經驗，專注雲端安全架構、滲透測試與事故應變。
%
\begin{itemize}
\item 熟悉零信任架構、SIEM、EDR 與 IAM，曾為金融及零售企業降低風險
\item 持有 CISSP、OSCP 及 CKS，擅長跨部門協作與合規審計
\item 以務實態度和清晰溝通推動安全左移，提升開發團隊安全意識
\end{itemize}}
\section{教育背景}

\cventry{2012年9月–2016年6月}
        {學士，電腦科學，成績：3.7}
        {香港科技大學}
        {\url{https://www.ust.hk}}
        {}
        {\begin{itemize}
\item 主修電腦科學，專注網絡與資訊安全
\item 畢業專題研究入侵偵測系統與機器學習應用
\item 曾獲院長嘉許名單
\end{itemize}
\textbf{課程}：資料結構與演算法、作業系統、電腦網絡、資料庫系統、資訊安全導論、密碼學基礎、雲端運算}
\section{工作經歷}

\cventry{2021年6月至今}
        {高級網絡安全工程師}
        {恒信金融科技有限公司}
        {\url{https://www.hengxin-fintech.com}}
        {}
        {\begin{itemize}
\item 設計及推行雲端安全架構，覆蓋 AWS 與 Azure，落實零信任存取及最少權限原則
\item 主導 SIEM 與 EDR 規則優化，將告警誤報率降低 35\%，平均事故回應時間縮短至 20 分鐘
\item 執行內部滲透測試與紅隊演練，協助修補高風險漏洞並建立修復追蹤流程
\item 與合規團隊完成 ISO 27001 與 PCI DSS 審計，維持零重大不符合項
\item 培訓開發團隊推行安全左移，將 SAST/DAST 納入 CI/CD 流程
\end{itemize}
\textbf{關鍵字}：雲端安全、零信任、事故應變、SIEM、安全左移}

\cventry{2018年3月–2021年5月}
        {網絡安全工程師}
        {駿達零售集團}
        {\url{https://www.chundaretail.com}}
        {}
        {\begin{itemize}
\item 負責零售支付系統的安全評估與監控，保護每日逾 50 萬筆交易
\item 部署及管理防火牆、WAF 與 IDS/IPS，定期進行規則檢視與調整
\item 主導弱點管理計劃，協調跨部門修補漏洞，季度高風險漏洞清零
\item 參與 ISO 27001 認證準備，撰寫風險評估與安全政策文件
\end{itemize}
\textbf{關鍵字}：滲透測試、弱點管理、ISO 27001、防火牆、WAF}

\cventry{2016年7月–2018年2月}
        {資訊安全分析師}
        {安網顧問有限公司}
        {\url{https://www.anwang-consulting.com}}
        {}
        {\begin{itemize}
\item 於 SOC 監控安全事件，進行初步分析、升級與事故記錄
\item 使用 Nessus、Burp Suite 等工具執行弱點掃描及滲透測試
\item 協助客戶完成安全基線檢查與合規差距分析
\item 撰寫事故報告與修復建議，向管理層匯報風險狀況
\end{itemize}
\textbf{關鍵字}：SOC、弱點掃描、風險評估、合規審計}
\section{語言}

\cvline{粵語}{母語或雙語水平 \hfill \textbf{關鍵字}：母語}
\cvline{英語}{完全專業水平 \hfill \textbf{關鍵字}：IELTS 7.5、商務溝通}
\cvline{普通話}{有限工作水平 \hfill \textbf{關鍵字}：日常會話}
\section{專業技能}

\cvline{網絡安全}{專家 \hfill \textbf{關鍵字}：滲透測試、紅隊演練、事故應變、弱點管理、風險評估}
\cvline{雲端與基礎架構}{高級 \hfill \textbf{關鍵字}：AWS、Azure、Kubernetes、Docker、Terraform}
\cvline{安全工具}{高級 \hfill \textbf{關鍵字}：Splunk、CrowdStrike、Burp Suite、Nessus、Wireshark}
\cvline{程式與自動化}{中級 \hfill \textbf{關鍵字}：Python、Bash、PowerShell、Git、Ansible}
\cvline{合規與治理}{高級 \hfill \textbf{關鍵字}：ISO 27001、PCI DSS、NIST CSF、零信任、安全政策}
\section{榮譽}

\cventry{2023年3月}
        {香港金融科技協會}
        {網絡安全卓越獎}
        {}
        {}
        {表揚在金融雲端安全轉型及事故應變流程優化上的突出貢獻。}

\cventry{2020年11月}
        {駿達零售集團}
        {年度最佳安全項目}
        {}
        {}
        {表彰在支付系統弱點管理及 ISO 27001 認證準備中的卓越表現。}
\section{證書}

\cventry{2022年8月}
        {ISC2}
        {CISSP}
        {\url{https://www.isc2.org/certifications/cissp}}
        {}
        {}

\cventry{2021年1月}
        {Offensive Security}
        {OSCP}
        {\url{https://www.offsec.com/courses/pen-200/}}
        {}
        {}

\cventry{2023年5月}
        {Linux Foundation}
        {CKS}
        {\url{https://training.linuxfoundation.org/certification/certified-kubernetes-security-specialist/}}
        {}
        {}

\cventry{2021年9月}
        {Amazon Web Services}
        {AWS Certified Security – Specialty}
        {\url{https://aws.amazon.com/certification/certified-security-specialty/}}
        {}
        {}
\section{出版刊物}

\cventry{2023年6月}
        {雲端原生環境的零信任存取實踐}
        {香港資訊安全期刊}
        {\url{https://www.hksecurityjournal.com/articles/zero-trust-cloud}}
        {}
        {\begin{itemize}
\item 探討零信任架構在混合雲環境的落地方法
\item 分享身份治理、微分段與持續驗證的實務經驗
\end{itemize}}
\section{項目}

\cventry{2022年1月–2023年12月}
        {為跨國金融機構設計及實施零信任安全架構}
        {企業零信任轉型計劃}
        {\url{https://www.chanhk-security.com/projects/zero-trust}}
        {}
        {\begin{itemize}
\item 重新設計身份驗證與存取控制流程，覆蓋 3,000 名員工及 200 個應用系統
\item 導入微分段及持續風險評估，降低橫向移動風險
\item 與 IT 及合規團隊協作，確保符合監管要求
\end{itemize}
\textbf{關鍵字}：零信任、身份治理、微分段、合規}

\cventry{2021年3月–2021年10月}
        {整合掃描、風險評分與修復追蹤的自動化平台}
        {自動化弱點管理平台}
        {\url{https://www.chanhk-security.com/projects/vuln-management}}
        {}
        {\begin{itemize}
\item 以 Python 串接 Nessus 與 Jira，自動建立修復工單
\item 建立風險導向評分模型，協助團隊優先處理高風險漏洞
\item 將平均修復週期由 45 天縮短至 18 天
\end{itemize}
\textbf{關鍵字}：Python、Nessus、Jira、自動化、風險評分}
\section{興趣愛好}

\cvline{網絡安全研究}{漏洞分析、威脅情報、CTF}
\cvline{戶外活動}{行山、單車、攝影}
\section{志願服務}

\cventry{2019年9月–2023年7月}
        {社區導師}
        {香港網絡安全意識協會}
        {\url{https://www.hkcyberawareness.org}}
        {}
        {\begin{itemize}
\item 為中學生及長者舉辦網絡安全講座，提升防騙意識
\item 協助設計釣魚郵件辨識工作坊教材
\item 參與年度網絡安全嘉年華的義工統籌
\end{itemize}}
    
\end{document}